\chapter{Discussion}
\label{cha:discussion}

On the circle \(S^1\), any target function can be decomposed into sine and
cosine waves. Low frequencies vary slowly around the circle, while high
frequencies oscillate more rapidly. The central question in this thesis is
whether training treats these frequency components separately, and whether the
preference for low frequencies survives once the ideal assumptions are relaxed.

In the infinite-width, infinite-data reference setting, the answer is explicit.
The NTK acts diagonally on the Fourier basis, meaning that each frequency
component of the residual evolves independently. The corresponding eigenvalue is
the learning rate of that frequency. Since the low-frequency eigenvalues are
larger, the low-frequency residual components decay faster. This is the
fixed-kernel spectral-bias mechanism in the setting studied here
\citep{jacot2018ntk,lee2020wide,rahaman2019spectralbiasneuralnetworks,
	basri2019convergence,bietti2019inductivebiasneuraltangent,
	cao2020understandingspectralbiasdeep}.

The continuum calculation also shows that the trainable hidden bias matters.
With the hidden-bias contribution, odd frequencies \(k\ge 3\) have nonzero
eigenvalues. Without this contribution, those frequencies are null directions of
the fixed-kernel dynamics. Thus the architecture does not only determine how
fast different frequencies are learned; it can also determine whether some
frequencies are learned at all in the fixed-kernel model. This agrees with
earlier harmonic analyses of two-layer ReLU networks, where bias-free models do
not express odd harmonics beyond the first frequency
\citep{basri2019convergence,basri2020frequencybiasneuralnetworks}.

\section{Answer to Q1: finite sampling}
\label{sec:discussion-q1}

The first research question asked what happens when the uniform continuum
measure on \(S^1\) is replaced by finitely many sampled input points.

The short answer is that exact frequency-wise learning is lost, but an
approximate version survives on fixed low-frequency blocks. After sampling, the
sine and cosine functions are only seen through their values on the sampled
points. They are no longer exactly orthonormal, and they are no longer exact
eigenvectors of the sampled NTK operator. However, the concentration bounds show
that these errors shrink as the number of samples increases. At fixed
confidence, the typical finite-sample error is of order \(n^{-1/2}\), up to
logarithmic factors.

This explains the sampled residual dynamics. Early in training, the retained
Fourier components decay at approximately the rates predicted by the corrected
finite-sample eigenvalues. Later, some components have already become very
small, so the remaining finite-sample error becomes visible relative to the
component being plotted. The late-time deviations therefore do not contradict
the Fourier prediction. They show the difference between exact diagonal
dynamics and an approximate sampled version of it.

The preconditioning experiment provides another check of this interpretation.
If most of the rate separation is caused by the eigenvalues, then rescaling by
those eigenvalues should make the retained components decay on more similar time
scales. This is what we observe. The experiment is related to work on modifying
kernel spectra through preconditioned dynamics
\citep{geifman2024controlling}, but its role here is mainly explanatory: it
tests whether the observed rate separation is mostly caused by the predicted
Fourier eigenvalues.

In more general data settings, one should not expect the Fourier basis itself to
remain the right reference. Fourier modes are natural here because the domain is
the circle and the data distribution is uniform. For non-uniform data or more
general domains, the analogous objects are the dominant eigenfunctions of the
kernel operator associated with the data distribution. The broader lesson is
therefore that finite samples should preserve the dominant population
eigendirections only approximately, with the largest visible errors appearing in
weaker or already-decayed components
\citep{bietti2019inductivebiasneuraltangent,
	basri2020frequencybiasneuralnetworks,cao2020understandingspectralbiasdeep}.

\section{Answer to Q2: finite width}
\label{sec:discussion-q2}

The second research question asked what happens when the network has finite
width. Here there are two cases: the tangent kernel can be frozen at
initialization, or it can evolve during training.

\subsection{Frozen tangent kernel}
\label{subsec:discussion-q2-frozen}

For the frozen tangent kernel, finite width behaves like a random perturbation
of the infinite-width model. On the retained low-frequency subspace
\(\mathcal H_K\), the projected finite-width operator concentrates around the
infinite-width operator at rate \(m^{-1/2}\), up to logarithmic factors. The
eigenspace plots show the same effect geometrically: as width increases, the
finite-width eigenspaces become closer to the Fourier frequency subspaces.

Thus, if the network is wide and the kernel is frozen, the infinite-width
Fourier picture remains a good approximation. The approximation is not equally
strong for all frequencies. Low frequencies have larger eigenvalues, so a small
perturbation is less important for their dynamics. Higher frequencies have
smaller eigenvalues and smaller gaps between nearby eigenvalues, so the same
perturbation can be more visible. This is consistent with the broader
spectral-bias literature, where high-frequency components are learned later and
are more sensitive to data distribution and model details
\citep{rahaman2019spectralbiasneuralnetworks,basri2019convergence,
	basri2020frequencybiasneuralnetworks}.

\subsection{Evolving tangent kernel}
\label{subsec:discussion-q2-evolving}

The evolving-kernel case is different. Here the kernel is not fixed during
training, so the fixed-kernel explanation cannot be the whole story. The results
are empirical rather than a theorem.

At small widths, allowing the kernel to evolve improves the final loss. At
larger widths, this benefit becomes smaller and eventually disappears. The
measurements suggest that the improvement at small width is not caused by better
alignment with the Fourier subspaces. Several higher-frequency subspaces become
less aligned during training. Instead, the main visible change is that the
evolving kernel increases operator strength in the lowest Fourier blocks,
especially \(k=0\) and \(k=1\), with a smaller increase for \(k=2\). This
matches the residual plots, where the strongest improvements occur in low and
intermediate target components.

The amplitude-variant experiments make this interpretation less dependent on
the particular target used in the main experiment. The low-frequency strength
increase still appears when all active non-constant modes have equal amplitudes,
and also when the larger amplitudes are assigned to higher frequencies. This
suggests that the effect is not simply caused by the original target putting
larger coefficients on low frequencies.

In more realistic networks, this distinction is important. A frozen-kernel
analysis can explain part of training when the network is wide enough and the
kernel changes little. But finite-width training can also change the kernel
itself, and this can change which directions are learned quickly. The broader
lesson is not that feature learning must always favor Fourier low frequencies.
Rather, feature learning can change the effective spectrum of the model, so a
purely fixed-kernel explanation may miss part of the training dynamics
\citep{bordelon2023dynamicsfinitewidthkernel,
	vyas2022limitationsntkunderstandinggeneralization}.

\section{Overall answer to the main research question}
\label{sec:discussion-main-question}

The main research question asked whether the simple frequency-wise learning
picture from the infinite-width, infinite-data NTK survives under more realistic
conditions.

For finite sampling, it survives approximately on fixed low-frequency blocks.
The sampled operator is not exactly diagonal in the Fourier basis, but the error
decreases with the number of samples, and the corrected eigenvalues explain the
early training dynamics.

For frozen finite width, it also survives approximately. The finite-width kernel
is random at initialization, but its projection onto \(\mathcal H_K\)
concentrates around the infinite-width operator as the width grows.

For evolving finite width, the fixed-kernel picture becomes incomplete. The
Fourier subspaces remain useful for measuring what happens, but the kernel
itself changes during training. In the experiments here, the most visible change
is an increase in low-frequency operator strength at small width.

The main conclusion is that finite sampling and frozen finite width behave like
controlled perturbations of the ideal NTK model on fixed low-frequency blocks.
Kernel evolution adds a further effect: training can reshape the operator
itself, and this can change the frequency preference of the dynamics.

\section{Limitations}
\label{sec:discussion-limitations}

\subsection{Scope of the concentration results}
\label{subsec:discussion-limitations-concentration}

The concentration results should be read as fixed-block perturbation results.
They control a prescribed low-frequency space \(\mathcal H_K\) with fixed \(K\).
They do not give a uniform statement over all Fourier modes, and the bounds are
not designed to remain sharp as \(K\) grows with \(n\) or \(m\). This is important
because the high-frequency eigenvalues and eigengaps become small, so controlling
higher and higher frequencies would require increasingly accurate operator
approximations.

The numerical constants are also conservative. In the finite-sample theorem, the
proof controls entries of the relevant matrices and then applies union bounds
over the retained block. The bound for \(H-\Lambda^{(n)}\) uses a
bounded-differences argument, and the bound for
\(A\Phi-\Phi\Lambda^{(n)}\) uses a conditional Bernstein estimate followed by a
union bound over sample-mode pairs. In the finite-width theorem, each projected
entry is controlled using a coarse tail envelope before another union bound is
applied. These steps are sufficient to prove concentration, but they discard
matrix structure and variance information. This explains why the theoretical
envelopes in Figures~\ref{fig:results-finite-sample-diagnostics} and
\ref{fig:results-finite-width-concentration} are much larger than the observed
errors.

In the finite-width concentration plot, the theorem contains an unspecified
universal constant \(c>0\). This constant comes from Bernstein's inequality for
averages of centered sub-exponential random variables
\citep{Vershynin_2018a}. In explicit proofs of such inequalities, constants of
order \(1/4\) often appear after converting a \(\psi_1\)-norm bound into a
moment generating function bound. The plotted choice \(c=0.2\) is a conservative
representative value, not an optimized constant for this problem. Changing this
constant shifts the envelope but does not change the predicted dependence on
\(m\).

A sharper analysis would need to use the structure of the objects being
controlled. For finite sampling, \(H\) is a bounded second-order statistic, so
concentration tools for U- or V-statistics may give tighter estimates. For both
finite sampling and finite width, direct matrix concentration in operator norm
would be preferable to entrywise control followed by norm conversion.

\subsection{Limits of the dynamics claims}
\label{subsec:discussion-limitations-dynamics}

The finite-sample dynamics experiments support the diagonal prediction, but the
theory does not give a complete time-uniform residual bound. The experiments
show that the predicted eigenvalues explain the dominant early-time behavior.
At later times, some residual components have already become very small, so even
a small difference between \(A\Phi\) and \(\Phi\Lambda^{(n)}\) can become visible
in relative terms. The current bounds identify the source of this effect, but
they do not give a sharp description of the full residual trajectory for all
times.

The preconditioning experiment is also a diagnostic rather than a proposed
training method. The theory preconditioner tests whether the observed rate
separation is mainly explained by the diagonal eigenvalue scaling. The empirical
preconditioner tests how much is gained by using the actual sampled
low-frequency block. These experiments support the interpretation of the
sampled dynamics, but they do not establish a general preconditioning algorithm
or a generalization guarantee.

\subsection{Limits of the finite-width alignment diagnostics}
\label{subsec:discussion-limitations-alignment}

The projector-error and principal-angle plots are empirical diagnostics of the
frozen finite-width operator. They compare Fourier frequency subspaces with
matched empirical eigenspaces of a discretized operator. This is useful for
checking whether the eigenspaces become more Fourier-like as width increases,
but it is not a separate theorem.

The matching procedure also introduces a practical limitation. For each
frequency \(k\), the empirical eigenspace is chosen by overlap with
\(\mathcal F_k\). This makes the comparison interpretable, but it also means
that the diagnostic depends on the selected discretization grid and on the
eigenspace-matching rule. The most useful conclusion from these plots is the trend
of errors in frequencies as width increases.

\subsection{Frozen versus evolving finite-width kernels}
\label{subsec:discussion-limitations-evolving}

The finite-width theorem concerns only the tangent operator at initialization.
It does not prove a bound for the evolving tangent kernel during training. The
frozen result therefore controls one source of finite-width error, namely random
initialization of the tangent kernel, but not the time-dependent change of the
kernel along the training trajectory.

The evolving-kernel section provides an empirical comparison between frozen and
evolving dynamics. It tracks the loss, mode-wise residual decay,
Fourier-subspace alignment, spectral strength, and block drift. These
diagnostics show a coherent width-dependent pattern, but they do not amount to a
time-dependent concentration theorem. A full theory would need to control the
evolving block \(C_t^{(K)}\), its drift from \(C_0^{(K)}\), and its action on the
current residual.

There is also no joint theorem combining finite sampling, finite width, and
training time. The finite-sample analysis studies the infinite-width sampled
operator. The finite-width theorem studies the continuum frozen operator. In
actual finite training runs, sampling error, finite-width initialization error,
and kernel evolution occur together.

\subsection{Model assumptions}
\label{subsec:discussion-limitations-model}

The analysis relies on the special structure of \(S^1\) with the uniform
measure. In this setting, the limiting kernel is rotation-invariant, so the
corresponding operator is a convolution operator and the Fourier basis
diagonalizes it exactly. Related harmonic descriptions of neural tangent kernels
on the circle or sphere appear in
\citet{basri2019convergence,basri2020frequencybiasneuralnetworks,bietti2019inductivebiasneuraltangent,cao2020understandingspectralbiasdeep}.
This structure is no longer automatic once the input distribution is changed.
For example, under a non-uniform density on the circle, the integral operator is
weighted by the density and is no longer diagonal in the global Fourier basis;
the relevant eigenfunctions become distribution-dependent, as studied by
\citet{basri2020frequencybiasneuralnetworks}. If the input domain is changed,
the appropriate basis also changes and may not be available in closed form.

The architecture and scaling assumptions are also part of the result. The thesis
uses a two-layer ReLU network in NTK scaling. The ReLU activation is important
because its Gaussian averages have explicit arc-cosine formulas
\citep{cho2009kernel}, which makes the kernel and its Fourier coefficients
computable in closed form. The \(1/\sqrt m\) scaling is important because it
places the model in the NTK or lazy-training regime, where the infinite-width
limit is described by a fixed kernel
\citep{jacot2018ntk,lee2020wide,golikov2022neuraltangentkernelsurvey}.
Changing the activation, adding or removing bias terms, changing the
parameterization, or moving to a feature-learning scaling changes the limiting
kernel and therefore changes the spectrum that drives the dynamics. Such
feature-learning regimes are better described by mean-field or related
infinite-width limits
\citep{Mei2018meanfieldview,chizat2018globalconvergencegradientdescent,sirignano2019meanfieldanalysisneural,Rotskoff_2022}.
Thus the Fourier conclusions should be understood as conclusions about this
explicit NTK model, not as architecture-independent statements about neural
networks in general.
